\documentclass[12pt,a4paper]{article}
\usepackage[utf8]{inputenc}
\usepackage[T1]{fontenc}
\usepackage{geometry}
\geometry{margin=1in, headheight=15pt}
\usepackage{setspace}
\usepackage{graphicx}
\usepackage{booktabs}
\usepackage{float}
\usepackage{amsmath}
\usepackage{tabularx}
\usepackage{longtable}
\usepackage{array}
\usepackage[svgnames,table]{xcolor}
\usepackage{titlesec}
\usepackage{fancyhdr}
\usepackage{microtype}
\usepackage{mathpazo}
\usepackage{caption}
\usepackage{tcolorbox}
\usepackage{hyperref}

\definecolor{primaryblue}{RGB}{20, 50, 120}
\definecolor{linkblue}{RGB}{0, 80, 160}

\hypersetup{
    colorlinks=true,
    linkcolor=primaryblue,
    filecolor=magenta,
    urlcolor=linkblue,
    citecolor=primaryblue,
    pdftitle={Sentiment Analysis for E-commerce Reviews},
    pdfauthor={Team SkyDom+}
}

\titleformat{\section}
  {\normalfont\Large\bfseries\color{primaryblue}}{\thesection}{1em}{}[{\color{primaryblue}\titlerule[0.8pt]}]
\titleformat{\subsection}
  {\normalfont\large\bfseries\color{primaryblue!85}}{\thesubsection}{1em}{}
\titleformat{\subsubsection}
  {\normalfont\normalsize\bfseries\color{primaryblue!85}}{\thesubsubsection}{1em}{}

\pagestyle{fancy}
\fancyhf{}
\lhead{\textcolor{gray}{\small NLP Project Report: Sentiment Analysis}}
\rhead{\textcolor{gray}{\small Team SkyDom+}}
\cfoot{\thepage}
\renewcommand{\headrulewidth}{0.4pt}
\renewcommand{\footrulewidth}{0.4pt}

\captionsetup{labelfont={bf,color=primaryblue}, textfont=it, margin=10pt}

\onehalfspacing
\setcounter{tocdepth}{2}
\renewcommand{\arraystretch}{1.25}

\begin{document}

\begin{titlepage}
\centering
\vspace*{0.1cm}
\includegraphics[width=0.7\textwidth]{\detokenize{D:/Code/NLP_code/p01_chart1.png}}\par
\vspace{1.0cm}
{\LARGE\bfseries\color{primaryblue} Sentiment Analysis for E-commerce Reviews\par}
\vspace{0.3cm}
{\Large\color{primaryblue!80} using Machine Learning Techniques.\par}
\vspace{0.8cm}
{\Large\bfseries Case Study Report\par}
\vspace{0.8cm}
{\large\bfseries Team: SkyDom+\par}
\vspace{0.3cm}
{\large Group members:\par}
\vspace{0.3cm}
\begin{center}
\large
\renewcommand{\arraystretch}{1.2}
\begin{tabular}{>{\raggedright\arraybackslash}p{6.0cm} >{\raggedright\arraybackslash}p{3.0cm}}
\toprule
\textbf{Name} & \textbf{ID} \\
\midrule
Nguyen Nhu Tien Tuan & 22BA13314 \\
Tran Hien Chuong & 22BA13056 \\
Do Gia Bao & 22BA13042 \\
Bui The Anh & 22BA13007 \\
Le Duc An & 22BA13002 \\
Pham Quang Tung & 22BA13320 \\
Khuong Bao Duc & 22BA13077 \\
\bottomrule
\end{tabular}
\end{center}
\vfill
{\large\bfseries March 12, 2026\par}
\end{titlepage}

\pagenumbering{roman}

\newpage
\renewcommand{\abstractname}{\Large\bfseries\color{primaryblue} Abstract}
\begin{abstract}
This report presents a consolidated NLP study in formal academic structure. The project addresses three tasks: binary sentiment classification with uncertainty control, multi-label issue extraction, and transformer-based extension. The source corpus contains 152,410 reviews; the issue annotation set contains 19,000 rows. In classic sentiment, variant V6 (word 1--2 + char 3--5 + negation, class weight w10) reaches test \texttt{precision\_0=0.677}, \texttt{recall\_0=0.932}, and \texttt{F2\_0=0.867}, with selective coverage 0.885 under threshold band 0.40/0.60. The transformer sentiment extension reaches \texttt{accuracy=0.966}, \texttt{precision\_0=0.799}, \texttt{recall\_0=0.853}, and \texttt{F2\_0=0.842}, with coverage 0.910. For issue extraction on the 3,800-row split, OVR Logistic reaches \texttt{micro-F1=0.780} and \texttt{macro-F1=0.387}, OVR LinearSVM reaches \texttt{micro-F1=0.771} and \texttt{macro-F1=0.405}, while dummy labelset baselines show high micro-F1 (\texttt{0.809}) but very low macro-F1 (\texttt{0.101}) due to majority-label dominance. The no-new-data upgrade package adds syllabus baselines, n-gram language modeling, evaluation-rigor artifacts, and transformer issue-routing experiments. All quantitative statements are sourced from saved experiment artifacts under \texttt{results/}.
\end{abstract}

\tableofcontents
\newpage
\pagenumbering{arabic}

\section{Introduction}
User reviews are short, noisy, and often mix praise and complaints in the same sentence. This project treats review understanding as a complete NLP pipeline problem rather than a single classifier benchmark \cite{jurafsky2026,turney2002}. The system is designed for practical support triage where missing negative feedback has higher operational risk than over-flagging.

\subsection{Project Goals}
The first goal is to build a sentiment analysis pipeline that can distinguish negative and positive reviews while handling uncertain cases more carefully than a standard hard classifier. The second goal is to move beyond coarse sentiment and extract concrete issue categories that can support downstream service action. The third goal is to compare classic sparse NLP with transformer-based contextual modeling under the same review domain.

\subsection{Target Users and Customers}
The primary target users are customer-support teams, operations teams, and product analysts working with large volumes of e-commerce feedback. A secondary target group is instructors and students in an NLP course setting, where the project can be used as a case study of how text classification, uncertainty handling, multi-label learning, and transformer extensions connect within one applied problem.

From an application perspective, the problem carries clear operational pain points. First, manual reading at scale is expensive and slow, so urgent negative cases can be buried among large volumes of positive comments. Second, a coarse positive/negative decision is not enough for action: support teams still need explicit issue types such as delivery, redemption, refund, or fraud signals. Third, short and ambiguous reviews create decision risk, where forcing hard labels can increase both false alarms and missed complaints.

These constraints motivate a design that goes beyond single-model accuracy. The project emphasizes recall-sensitive evaluation for risky negative cases, explicit issue extraction for downstream handling, and uncertainty-aware outputs for low-information texts. This framing links NLP model choices directly to practical service quality and response time.

The scope is constrained to no-new-data upgrades: modeling, evaluation, calibration, and analysis improvements are performed on existing data only.

\section{NLP Problem Definition}
\subsection{Task A: Sentiment Understanding}
Task A is binary sentiment classification from review text. Ratings 1--2 are mapped to class 0 (negative), ratings 4--5 to class 1 (positive), and rating 3 is reserved as uncertain analysis data. Due to severe imbalance (about 9.5:1 positive to negative), model selection follows a negative-first ordering: maximize \texttt{recall\_0}, then \texttt{F2\_0}, then \texttt{precision\_0}.

\subsection{Task B: Issue Extraction}
Task B maps review text to a multi-label issue vector with nine labels (delivery, redemption, product quality, customer service, refund, usability, value/price, fraud/scam, other). This converts unstructured text into actionable support categories \cite{tsoumakas2011,zhang2014}.

\subsection{Task C: Deep Context Extension}
Task C extends the classic pipeline with transformer fine-tuning for deeper contextual modeling, hard-case comparison, and threshold-behavior analysis \cite{vaswani2017,devlin2019}.

\section{Dataset and Language Characteristics}
\subsection{Data Source and Scale}
The source file contains 152,410 reviews. Rating distribution is shown in Table \ref{tab:rating_distribution}.

\begin{table}[H]
\centering
\caption{Rating distribution in source data}
\label{tab:rating_distribution}
\begin{tabular}{lr}
\toprule
Rating & Count \\
\midrule
1 & 12,326 \\
2 & 1,873 \\
3 & 3,271 \\
4 & 6,692 \\
5 & 128,248 \\
\bottomrule
\end{tabular}
\end{table}

Table \ref{tab:rating_distribution} confirms an extreme long-tail sentiment distribution in raw e-commerce reviews. Rating 5 alone accounts for roughly 84.1\% of all rows, while rating 1 and 2 together account for about 9.3\%. This imbalance explains why accuracy-only evaluation is insufficient and why the project prioritizes negative-class recall and F2. It also justifies the decision to keep rating 3 as an uncertainty set instead of forcing it into hard supervision.

Strong-label sentiment rows (ratings 1,2,4,5) total 149,139, with 14,199 negatives and 134,940 positives.

\subsection{Observed Language Noise}
The corpus includes abbreviations (for example \textit{gr8}, \textit{thx}, \textit{idk}), short fragments, typos, and mixed-polarity constructions such as \textit{good but late delivery}. Negation and contrast markers are frequent and motivate dedicated handling in preprocessing and feature design.

\subsection{Issue Label Data}
The merged issue dataset contains 19,000 annotated rows. Label frequency is shown in Table \ref{tab:issue_freq}.

\begin{table}[H]
\centering
\caption{Issue label frequency in manual annotation set}
\label{tab:issue_freq}
\begin{tabular}{lrr}
\toprule
Label & Count & Share \\
\midrule
delivery\_shipping & 972 & 0.0512 \\
redemption\_activation & 284 & 0.0149 \\
product\_quality & 1,251 & 0.0658 \\
customer\_service & 182 & 0.0096 \\
refund\_return & 93 & 0.0049 \\
usability & 873 & 0.0459 \\
value\_price & 558 & 0.0294 \\
fraud\_scam & 144 & 0.0076 \\
other & 15,960 & 0.8400 \\
\bottomrule
\end{tabular}
\end{table}

The issue distribution in Table \ref{tab:issue_freq} is even more skewed than sentiment labels: \texttt{other} dominates at 84.0\%, and several actionable labels remain below 1.0\% prevalence. This creates a high risk that naive models appear strong on aggregate metrics while missing operationally important minority issues. Therefore, macro-F1, per-label F1, and per-label threshold tuning are necessary to avoid majority-label bias in deployment decisions.

\section{Text Processing Design}
\subsection{Normalization}
The default cleaner applies lowercasing, URL removal, punctuation spacing, and whitespace normalization. The goal is to reduce sparsity while preserving sentiment-bearing tokens.

\subsection{Negation Handling}
The classic sentiment and issue pipelines support negation propagation with a local window (current issue config: window size 3). This is critical for polarity flips in short reviews.

\subsection{Contrast Handling}
Contrast cues such as \textit{but}, \textit{however}, and clause-level polarity shifts are handled indirectly through n-gram features, including phrase patterns that encode polarity transitions.

\subsection{Character-level Modeling}
Character n-grams (3--5) are combined with word n-grams in selected variants to improve robustness to misspellings, abbreviations, and non-standard surface forms.

\section{Feature Engineering}
\subsection{Classic Sparse Representation}
The main sentiment representation is TF-IDF over word n-grams (1--2) and character n-grams (3--5), followed by Chi-square feature selection. The selected best variant (V6) uses \texttt{k=10000}.

For issue extraction, TF-IDF with optional character support and Chi-square filtering (selected \texttt{k=2000}) is used before one-vs-rest classification.

\subsection{Lexicon Feature Block}
The current default best models are data-driven sparse models. Lexicon-oriented analyses are retained as optional diagnostic components in the extension package rather than as mandatory blocks in the final deployment path.

\section{Modeling Approach}
\subsection{Sentiment Classifiers}
The classic sentiment backbone is Logistic Regression. Class-weight modes are searched across \texttt{none}, \texttt{balanced}, \texttt{w2}, \texttt{w5}, and \texttt{w10}. Variant search over preprocessing and feature toggles (V0--V7) determines the final configuration.

\subsection{Issue Extraction Model}
Issue extraction uses calibrated one-vs-rest Logistic Regression with per-label threshold tuning. One-vs-rest LinearSVM and a probability-level blend (\texttt{ovr\_blend\_lr\_svm}) are retained as strong classical comparators. Two dummy baselines (\texttt{dummy\_labelset\_majority}, \texttt{dummy\_label\_prior}) are also reported as sanity checks. Current train configuration includes per-label class-weight search and sigmoid calibration (CV=3).

\subsection{Transformer Extension}
The transformer extension uses HuggingFace sequence-classification fine-tuning (default base family: DistilBERT in this setup). It is evaluated with the same threshold-band policy to enable direct comparison with classic selective behavior.

\section{Experimental Setup}
\subsection{Split and Protocol}
Sentiment training uses strong-label data (ratings 1,2,4,5) with held-out validation and test sets; rating 3 is kept for uncertainty analysis. Issue extraction uses a fixed random split of 13,300/1,900/3,800 (train/val/test). The issue-transformer extension reports on 3,000/1,000/2,000.

\subsection{Main Metrics}
For sentiment, main metrics are \texttt{accuracy}, \texttt{F1}, \texttt{precision\_0}, \texttt{recall\_0}, \texttt{F1\_0}, and \texttt{F2\_0}. For issue extraction, metrics are micro-F1, macro-F1, subset accuracy, hamming loss, and per-label F1.

\subsection{Uncertainty Rule}
Selective classification uses a threshold band \([t_{low}, t_{high}]\). Samples with probability inside the band are marked uncertain; remaining samples are covered decisions. Reported selective metrics include coverage, selective precision\_0, selective recall\_0, selective F2\_0, and covered false-negative rate \cite{geifman2017}.

\section{No-New-Data Syllabus Upgrade Package}
The upgrade package extends the project without collecting new data. Major components in the current experimental package are:
\begin{enumerate}
\item Issue Level-3 enhancements: per-label class-weight search, probability calibration, threshold stability diagnostics.
\item Issue Level-4 transformer baseline and hybrid routing with capped route rate.
\item Syllabus benchmark suite: Naive Bayes, Perceptron, LinearSVM, SGD log-loss, semantic SVD baselines, FFNN MLP on SVD features.
\item N-gram language model module with add-k smoothing and perplexity reporting.
\item Evaluation-rigor package: bootstrap confidence intervals, pairwise significance tests, and error taxonomy.
\item Course-fit matrix generated directly from saved artifacts.
\end{enumerate}

\section{Results for Classic NLP Sentiment Model}
\subsection{Baseline Selection by K}
Table \ref{tab:chi2_k} shows best validation setting per Chi-square dimension under negative-first ordering.

\begin{table}[H]
\centering
\caption{Best sentiment baseline by Chi-square K (validation)}
\label{tab:chi2_k}
\begin{tabular}{lccccc}
\toprule
K & Class weight & Precision\_0 & Recall\_0 & F2\_0 & Accuracy \\
\midrule
1,000 & w10 & 0.555 & 0.910 & 0.806 & 0.922 \\
2,000 & w10 & 0.584 & 0.920 & 0.825 & 0.930 \\
5,000 & balanced & 0.606 & 0.935 & 0.843 & 0.936 \\
10,000 & balanced & 0.621 & 0.936 & 0.850 & 0.940 \\
\bottomrule
\end{tabular}
\end{table}

Table \ref{tab:chi2_k} shows a monotonic quality gain as feature dimension increases from 1,000 to 10,000. Precision\_0 improves from 0.555 to 0.621 and F2\_0 rises from 0.806 to 0.850, while recall\_0 stays high across all settings. This pattern indicates that larger sparse vocabularies mainly improve false-positive control on the negative class without sacrificing triage sensitivity.

\subsection{Variant Comparison}
Table \ref{tab:classic_variants} reports test performance for engineered variants. V6 is the selected main model.

\begin{table}[H]
\centering
\caption{Classic sentiment variant comparison on test}
\label{tab:classic_variants}
\begin{tabular}{lccccc}
\toprule
Variant & K & Class weight & Precision\_0 & Recall\_0 & F2\_0 \\
\midrule
V0 & 10,000 & balanced & 0.638 & 0.937 & 0.856 \\
V1 & 10,000 & balanced & 0.643 & 0.932 & 0.855 \\
V2 & 5,000 & w10 & 0.647 & 0.926 & 0.852 \\
V3 & 10,000 & balanced & 0.626 & 0.923 & 0.843 \\
V4 & 10,000 & balanced & 0.635 & 0.921 & 0.845 \\
V5 & 10,000 & balanced & 0.653 & 0.933 & 0.859 \\
V6 & 10,000 & w10 & 0.677 & 0.932 & 0.867 \\
V7 & 5,000 & w10 & 0.657 & 0.927 & 0.857 \\
\bottomrule
\end{tabular}
\end{table}

Compared with dummy sentiment baselines (\texttt{dummy\_most\_frequent}: \texttt{F2\_0=0.000}, \texttt{dummy\_stratified}: \texttt{F2\_0=0.092}), the selected V6 model (\texttt{F2\_0=0.867}) shows a large practical margin on the negative class objective.
Within the non-dummy family, V6 also gives the best balance: it achieves the highest precision\_0 (0.677) while retaining recall\_0 above 0.93, outperforming nearby variants V5 and V7 on F2\_0. This suggests that the joint use of word+character TF-IDF plus negation modeling contributes complementary signals rather than redundant complexity.

\subsection{Threshold and Uncertainty Behavior}
For V6, the recommended threshold band is 0.40/0.60. Table \ref{tab:classic_threshold} shows sensitivity to threshold movement.

\begin{table}[H]
\centering
\caption{Classic sentiment threshold sweep on test}
\label{tab:classic_threshold}
\begin{tabular}{lcccc}
\toprule
Thresholds & Coverage & Sel. precision\_0 & Sel. recall\_0 & Sel. F2\_0 \\
\midrule
0.40 / 0.60 & 0.885 & 0.730 & 0.956 & 0.900 \\
0.45 / 0.60 & 0.891 & 0.708 & 0.956 & 0.894 \\
0.50 / 0.65 & 0.886 & 0.681 & 0.963 & 0.889 \\
0.55 / 0.70 & 0.881 & 0.646 & 0.970 & 0.882 \\
\bottomrule
\end{tabular}
\end{table}

Table \ref{tab:classic_threshold} reveals a controlled precision-recall trade-off as thresholds move upward. Increasing the positive-decision band improves selective recall\_0 (0.956 to 0.970) but reduces selective precision\_0 (0.730 to 0.646), while coverage remains stable around 0.88--0.89. This stability is valuable for operations because policy adjustments can tune risk profile without large throughput disruption.

\section{Results for Transformer NLP Extension}
\subsection{Transformer Metrics}
Table \ref{tab:tf_metrics} summarizes transformer sentiment metrics.

\begin{table}[H]
\centering
\caption{Transformer sentiment performance}
\label{tab:tf_metrics}
\begin{tabular}{lcccccc}
\toprule
Split & Accuracy & F1 & Precision\_0 & Recall\_0 & F2\_0 & Coverage \\
\midrule
Validation & 0.966 & 0.981 & 0.793 & 0.870 & 0.854 & 0.906 \\
Test & 0.966 & 0.981 & 0.799 & 0.853 & 0.842 & 0.910 \\
\bottomrule
\end{tabular}
\end{table}

Table \ref{tab:tf_metrics} indicates strong generalization consistency: validation and test metrics are closely aligned, with no sign of major overfitting in one-epoch fine-tuning. Compared to the classic model, transformer improves negative precision materially (0.799 vs 0.677) but reduces negative recall (0.853 vs 0.932). This means transformer is better at avoiding false alarms, while classic V6 remains better at catching risky negative cases under the project’s triage-first objective.

\subsection{Transformer Threshold Sweep}
\begin{table}[H]
\centering
\caption{Transformer threshold sweep on test}
\label{tab:tf_threshold}
\begin{tabular}{lcccc}
\toprule
Thresholds & Coverage & Sel. precision\_0 & Sel. recall\_0 & Sel. F2\_0 \\
\midrule
0.40 / 0.60 & 0.910 & 0.831 & 0.846 & 0.843 \\
0.45 / 0.60 & 0.911 & 0.827 & 0.847 & 0.843 \\
0.50 / 0.65 & 0.912 & 0.821 & 0.850 & 0.844 \\
0.55 / 0.70 & 0.911 & 0.818 & 0.855 & 0.847 \\
\bottomrule
\end{tabular}
\end{table}

Table \ref{tab:tf_threshold} shows the transformer selective curve is relatively flat, with coverage near 0.91 across all tested bands and only small F2\_0 movement (0.843 to 0.847). This suggests the model is calibration-stable in the evaluated range and can be tuned conservatively without destabilizing decision volume.

\subsection{Hard-case Comparison}
Saved hard-case probes show complementary behavior. For \textit{not bad} and \textit{good but late delivery}, the classic model predicts negative while the transformer predicts positive. For very short tokens (\textit{gr8}, \textit{thx}, \textit{idk}), both systems route to uncertain due to short-text policy.

\section{Results for NLP Issue Extraction}
\subsection{Overall Multi-label Results}
Table \ref{tab:issue_overall} reports the current main pipeline results (test size 3,800).

\begin{table}[H]
\centering
\caption{Issue extraction overall metrics}
\label{tab:issue_overall}
\begin{tabular}{lcccc}
\toprule
Model & Micro-F1 & Macro-F1 & Subset accuracy & Hamming loss \\
\midrule
OVR Logistic (test) & 0.780 & 0.387 & 0.634 & 0.0585 \\
OVR LinearSVM (test) & 0.771 & 0.405 & 0.617 & 0.0623 \\
OVR Blend LR+SVM (test) & 0.768 & 0.391 & 0.626 & 0.0620 \\
Dummy labelset majority (test) & 0.809 & 0.101 & 0.811 & 0.0438 \\
Dummy label prior (test) & 0.809 & 0.101 & 0.811 & 0.0438 \\
\bottomrule
\end{tabular}
\end{table}

Dummy baselines rank high on micro-F1 and subset accuracy because they almost always predict the dominant \texttt{other} label; macro-F1 exposes that they fail on minority issue labels.
From Table \ref{tab:issue_overall}, OVR Logistic is preferred for balanced production utility: it leads micro-F1 and subset accuracy among non-dummy models, while LinearSVM offers slightly better macro-F1. The blend does not beat both parents, implying that simple averaging is not sufficient under the current label imbalance and would require stronger routing or calibration logic.

\subsection{Per-label Improvement}
Table \ref{tab:issue_delta} compares per-label F1 between SVM and Logistic on test.

\begin{table}[H]
\centering
\caption{Per-label F1 change (SVM minus Logistic)}
\label{tab:issue_delta}
\begin{tabular}{lccc}
\toprule
Label & Logistic F1 & SVM F1 & $\Delta$ \\
\midrule
delivery\_shipping & 0.442 & 0.446 & +0.004 \\
redemption\_activation & 0.298 & 0.287 & -0.010 \\
product\_quality & 0.307 & 0.282 & -0.025 \\
customer\_service & 0.225 & 0.168 & -0.057 \\
refund\_return & 0.435 & 0.387 & -0.048 \\
usability & 0.459 & 0.462 & +0.003 \\
value\_price & 0.258 & 0.313 & +0.056 \\
fraud\_scam & 0.150 & 0.391 & +0.241 \\
other & 0.911 & 0.911 & +0.001 \\
\bottomrule
\end{tabular}
\end{table}

Table \ref{tab:issue_delta} highlights where SVM helps versus hurts. The largest gain is on \texttt{fraud\_scam} (+0.241), followed by \texttt{value\_price} (+0.056), indicating better separation on specific sparse semantic cues. However, Logistic remains better on \texttt{customer\_service}, \texttt{refund\_return}, and \texttt{product\_quality}, which are critical support labels; this is why the main pipeline keeps Logistic as default and uses SVM as a comparator rather than replacement.

\subsection{Threshold Strategy in Multi-label Setting}
The tuned per-label thresholds for the main \texttt{ovr\_logreg} pipeline are shown in Table \ref{tab:issue_thresholds}. Lower thresholds are assigned to rare labels to improve recall, while \texttt{fraud\_scam} and \texttt{other} are kept more conservative.

\begin{table}[H]
\centering
\caption{Final per-label thresholds in issue pipeline}
\label{tab:issue_thresholds}
\begin{tabular}{lc}
\toprule
Label & Threshold \\
\midrule
delivery\_shipping & 0.25 \\
redemption\_activation & 0.10 \\
product\_quality & 0.10 \\
customer\_service & 0.10 \\
refund\_return & 0.15 \\
usability & 0.10 \\
value\_price & 0.15 \\
fraud\_scam & 0.50 \\
other & 0.30 \\
\bottomrule
\end{tabular}
\end{table}

Threshold values in Table \ref{tab:issue_thresholds} encode an explicit policy: rare but important labels are granted low activation cutoffs (0.10--0.15) to reduce missed incidents, while high-risk false-positive channels are constrained with higher cutoffs (\texttt{fraud\_scam}=0.50, \texttt{other}=0.30). This per-label asymmetry is more realistic than a global threshold and better aligned with downstream ticketing workflows.

\section{Results for Syllabus Upgrade Package}
\subsection{Classic Baselines for Core Course Algorithms}
Table \ref{tab:syllabus_baselines} summarizes no-new-data sentiment baselines.

\begin{table}[H]
\centering
\caption{Syllabus benchmark baselines on test}
\label{tab:syllabus_baselines}
\begin{tabular}{lccc}
\toprule
Model & Recall\_0 & Precision\_0 & F2\_0 \\
\midrule
logreg\_l2 & 0.913 & 0.703 & 0.862 \\
multinomial\_nb & 0.913 & 0.633 & 0.839 \\
complement\_nb & 0.952 & 0.510 & 0.811 \\
perceptron & 0.828 & 0.684 & 0.795 \\
linear\_svm & 0.847 & 0.789 & 0.835 \\
sgd\_log\_loss & 0.921 & 0.616 & 0.838 \\
svd\_semantic\_logreg & 0.920 & 0.501 & 0.788 \\
ffnn\_mlp\_svd & 0.730 & 0.799 & 0.742 \\
\bottomrule
\end{tabular}
\end{table}

Table \ref{tab:syllabus_baselines} provides a clear algorithmic trade-off map for defense: Complement NB maximizes recall\_0 (0.952) but with weak precision\_0 (0.510), while LinearSVM gives high precision\_0 (0.789) at lower recall\_0. The best negative-class balance under F2 is achieved by logistic regression (\texttt{F2\_0=0.862}), which supports its role as the primary classic baseline in this project.

\subsection{Selective Behavior in the Added Baselines}
\begin{table}[H]
\centering
\caption{Selective metrics for probabilistic syllabus baselines (test)}
\label{tab:syllabus_selective}
\begin{tabular}{lcccc}
\toprule
Model & Coverage & Sel. precision\_0 & Sel. recall\_0 & Sel. F2\_0 \\
\midrule
logreg\_l2\_selective & 0.887 & 0.763 & 0.946 & 0.903 \\
multinomial\_nb\_selective & 0.903 & 0.654 & 0.928 & 0.857 \\
complement\_nb\_selective & 0.900 & 0.545 & 0.963 & 0.835 \\
sgd\_log\_loss\_selective & 0.863 & 0.696 & 0.959 & 0.892 \\
\bottomrule
\end{tabular}
\end{table}

Selective results in Table \ref{tab:syllabus_selective} show that \texttt{logreg\_l2\_selective} remains strongest on Sel.\ F2\_0 (0.903), combining high Sel.\ recall\_0 (0.946) with solid precision\_0 (0.763). Complement NB still reaches the highest Sel.\ recall\_0 (0.963) but precision collapses, confirming it is suitable only when maximum capture is prioritized over review workload.

\subsection{N-Gram Language Model Results}
\begin{table}[H]
\centering
\caption{N-gram LM perplexity (add-k smoothing)}
\label{tab:ngram_ppl}
\begin{tabular}{lcc}
\toprule
Model & Validation perplexity & Test perplexity \\
\midrule
Unigram add-k & 376.51 & 378.02 \\
Bigram add-k & 358.59 & 363.08 \\
\bottomrule
\end{tabular}
\end{table}

Table \ref{tab:ngram_ppl} confirms expected language-model behavior: bigram consistently outperforms unigram on both validation and test perplexity. On test data, perplexity drops from 378.02 to 363.08 (about 3.95\% relative reduction), showing that even short review texts benefit from one-step context modeling. This provides concrete evidence for the course N-gram component beyond classification metrics.

\subsection{Updated Course-Fit Matrix}
Table \ref{tab:course_fit} reports the topic-coverage matrix generated from current artifacts. The mean coverage is \textbf{94.3\%}.

\begin{table}[H]
\centering
\caption{Course-fit matrix from current artifact set}
\label{tab:course_fit}
\begin{tabular}{lc}
\toprule
Topic & Coverage (\%) \\
\midrule
Probability & 100 \\
Regular Expressions & 50 \\
Evaluation Measures & 100 \\
N-Gram Language Models & 100 \\
Naive Bayes & 100 \\
Perceptron & 100 \\
Logistic Regression & 100 \\
Feed-forward Neural Networks & 100 \\
Pytorch Basic & 100 \\
Vector Semantics and Embeddings & 80 \\
Recurrent Neural Networks & 100 \\
Transformers and LLMs & 100 \\
Masked Language Models & 100 \\
Applications of LLMs & 90 \\
\bottomrule
\end{tabular}
\end{table}

The matrix in Table \ref{tab:course_fit} indicates strong breadth across core topics, with most items at 100\%. Remaining gaps are concentrated in \textit{Regular Expressions} (50\%) and partly in \textit{Vector Semantics and Embeddings} (80\%), so these are the highest-impact targets if the team wants to further improve academic completeness without changing dataset scope.

\section{Classic Versus Transformer Discussion}
On sentiment classification, transformer improves overall accuracy and negative precision (0.799 vs 0.677), while classic V6 remains stronger on negative recall (0.932 vs 0.853) and F2\_0 (0.867 vs 0.842). Under the same threshold band (0.40/0.60), transformer yields higher selective precision but lower selective recall than classic.

For issue extraction, Table \ref{tab:issue_fair} reports the dedicated fair-split benchmark. The classic issue model remains strongest; hybrid routing improves over pure transformer but still trails the classic setting.

\begin{table}[H]
\centering
\caption{Issue extraction fair-split comparison}
\label{tab:issue_fair}
\begin{tabular}{lcccc}
\toprule
Model & Micro-F1 & Macro-F1 & Subset accuracy & Hamming loss \\
\midrule
classic\_issue\_model & 0.792 & 0.528 & 0.656 & 0.0549 \\
hybrid\_route & 0.730 & 0.338 & 0.615 & 0.0740 \\
transformer\_multilabel & 0.646 & 0.188 & 0.551 & 0.1013 \\
\bottomrule
\end{tabular}
\end{table}

Relative to pure transformer on this fair split, the classic issue model gains +0.146 micro-F1 and +0.340 macro-F1.
Table \ref{tab:issue_fair} also shows that hybrid routing recovers part of the transformer gap but remains clearly below the classic model, especially on macro-F1 where minority-label handling is critical. This suggests transformer issue modeling is not yet data-efficient enough under current annotation volume and should be treated as an auxiliary component until label support increases or architecture is adapted.

\section{Error Analysis and Failure Modes}
The evaluation-rigor package records 1,454 sentiment misclassified rows (194 false negatives for class 0, 1,260 false positives for class 0). Dominant false-positive categories are negation patterns (49.8\%), other mixed contexts (32.4\%), and contrast patterns (10.4\%). Dominant false-negative categories are other contexts (52.1\%), short text (23.2\%), and negation patterns (14.4\%).

Pairwise significance tests against baseline comparators show positive deltas for V6 (for example recall\_0 +0.063 vs multinomial NB, +0.085 vs LinearSVM), but p-values remain above 0.05 in this run, so gains are treated as directional rather than statistically confirmed.

\section{Practical Value for NLP Course}
The project now provides a full, reproducible course artifact chain: cleaned datasets, stepwise experimental outputs, saved models, threshold analyses, uncertainty policy, transformer extension, and evaluation-rigor outputs. This supports both engineering deployment discussion and academic defense requirements.

The upgraded syllabus package strengthens mapping to core course topics while preserving the same data budget, making the project suitable for methodology-focused assessment. Based on measurable experimental artifacts, the latest rubric assessment estimates 18.04/20 (90.2\%) and syllabus-topic coverage 94.3\%.

\section{Limitations and Future Work}
Main limitations are class imbalance, high dominance of the \texttt{other} issue label, and weak transformer multi-label performance under current small-split fine-tuning. Error taxonomy also shows persistent sensitivity to negation and mixed-contrast sentences.

Future work should prioritize calibration-aware ensembling, better contextual augmentation without new labels, and lightweight sequence models for issue extraction that preserve interpretability.

\section{Conclusion}
This report summarizes model structure and quantitative findings in one consistent analysis. The main production model for sentiment is still the classic V6 variant because it optimizes negative recall/F2 under project policy. Transformer sentiment provides stronger precision and overall accuracy, and remains valuable for complementary analysis. For issue extraction, calibrated classic OVR models remain the strongest option in the current experiments. The no-new-data upgrade package substantially improves academic completeness, reproducibility, and course-topic coverage.

\clearpage
\section*{References}
\addcontentsline{toc}{section}{References}

\begin{thebibliography}{99}

\bibitem{jurafsky2026}
Dan Jurafsky and James H. Martin.
\textit{Speech and Language Processing} (3rd ed. draft).
2026.
\newline\texttt{https://web.stanford.edu/~jurafsky/slp3/}

\bibitem{turney2002}
Peter D. Turney.
Thumbs Up or Thumbs Down? Semantic Orientation Applied to Unsupervised Classification of Reviews.
In \textit{Proceedings of ACL}, 2002.
\newline\texttt{https://aclanthology.org/P02-1053/}

\bibitem{wang2012}
Sida Wang and Christopher D. Manning.
Baselines and Bigrams: Simple, Good Sentiment and Topic Classification.
In \textit{Proceedings of ACL}, 2012.
\newline\texttt{https://aclanthology.org/P12-2018/}

\bibitem{tsoumakas2011}
Grigorios Tsoumakas, Ioannis Katakis, and Ioannis Vlahavas.
Random k-Labelsets for multilabel classification.
\textit{IEEE Transactions on Knowledge and Data Engineering}, 23(7):1079--1089, 2011.
\newline\texttt{https://doi.org/10.1109/TKDE.2010.164}

\bibitem{zhang2014}
Min-Ling Zhang and Zhi-Hua Zhou.
A review on multi-label learning algorithms.
\textit{IEEE Transactions on Knowledge and Data Engineering}, 26(8):1819--1837, 2014.
\newline\texttt{https://doi.org/10.1109/TKDE.2013.39}

\bibitem{vaswani2017}
Ashish Vaswani, Noam Shazeer, Niki Parmar, Jakob Uszkoreit, Llion Jones, Aidan N. Gomez, Lukasz Kaiser, and Illia Polosukhin.
Attention Is All You Need.
In \textit{Advances in Neural Information Processing Systems 30}, 2017.
\newline\texttt{https://arxiv.org/abs/1706.03762}

\bibitem{devlin2019}
Jacob Devlin, Ming-Wei Chang, Kenton Lee, and Kristina Toutanova.
BERT: Pre-training of Deep Bidirectional Transformers for Language Understanding.
In \textit{Proceedings of NAACL-HLT}, 2019.
\newline\texttt{https://aclanthology.org/N19-1423/}

\bibitem{saito2015}
Takaya Saito and Marc Rehmsmeier.
The precision-recall plot is more informative than the ROC plot when evaluating binary classifiers on imbalanced datasets.
\textit{PLOS ONE}, 10(3):e0118432, 2015.
\newline\texttt{https://doi.org/10.1371/journal.pone.0118432}

\bibitem{geifman2017}
Yonatan Geifman and Ran El-Yaniv.
Selective Classification for Deep Neural Networks.
\textit{arXiv preprint arXiv:1705.08500}, 2017.
\newline\texttt{https://arxiv.org/abs/1705.08500}

\end{thebibliography}

\clearpage
\section*{Appendix A: Sentiment Figures}
\addcontentsline{toc}{section}{Appendix A: Sentiment Figures}

\begin{figure}[H]
\centering
\includegraphics[width=0.82\textwidth]{\detokenize{results/dm2_steps/01_rating_distribution.png}}
\caption{Rating distribution in source data}
\end{figure}

\begin{figure}[H]
\centering
\includegraphics[width=0.82\textwidth]{\detokenize{results/dm2_steps/06_f1_vs_k.png}}
\caption{Validation F1 by Chi-square top K}
\end{figure}

\begin{figure}[H]
\centering
\includegraphics[width=0.82\textwidth]{\detokenize{results/dm2_steps/06_recall0_vs_k.png}}
\caption{Validation recall for negative class by top K}
\end{figure}

\begin{figure}[H]
\centering
\includegraphics[width=0.82\textwidth]{\detokenize{results/dm2_steps/06b_recall0_by_variant.png}}
\caption{Test recall for negative class by classic variant}
\end{figure}

\begin{figure}[H]
\centering
\includegraphics[width=0.82\textwidth]{\detokenize{results/dm2_steps/06b_precision0_by_variant.png}}
\caption{Test precision for negative class by classic variant}
\end{figure}

\begin{figure}[H]
\centering
\includegraphics[width=0.82\textwidth]{\detokenize{results/dm2_steps/09_prob_hist_test_strong.png}}
\caption{Classic probability histogram on strong test set}
\end{figure}

\begin{figure}[H]
\centering
\includegraphics[width=0.82\textwidth]{\detokenize{results/dm2_steps/10_threshold_tradeoff.png}}
\caption{Classic threshold trade-off for coverage and recall}
\end{figure}

\begin{figure}[H]
\centering
\includegraphics[width=0.82\textwidth]{\detokenize{results/nlp_ext/nlp_confusion_matrix.png}}
\caption{Transformer confusion matrix on test set}
\end{figure}

\begin{figure}[H]
\centering
\includegraphics[width=0.82\textwidth]{\detokenize{results/nlp_ext/nlp_prob_hist_test.png}}
\caption{Transformer probability histogram on test set}
\end{figure}

\begin{figure}[H]
\centering
\includegraphics[width=0.82\textwidth]{\detokenize{results/nlp_ext/nlp_threshold_tradeoff.png}}
\caption{Transformer threshold trade-off for coverage and recall}
\end{figure}

\clearpage
\section*{Appendix B: Issue Extraction Figures}
\addcontentsline{toc}{section}{Appendix B: Issue Extraction Figures}

\begin{figure}[H]
\centering
\includegraphics[width=0.82\textwidth]{\detokenize{results/issue_steps/02_label_frequency.png}}
\caption{Issue label frequency in annotation set}
\end{figure}

\begin{figure}[H]
\centering
\includegraphics[width=0.82\textwidth]{\detokenize{results/issue_steps/02_per_label_f1.png}}
\caption{Per-label F1 for OVR Logistic issue model}
\end{figure}

\begin{figure}[H]
\centering
\includegraphics[width=0.82\textwidth]{\detokenize{results/issue_steps_char_demo/02_per_label_f1.png}}
\caption{Per-label F1 for word+char issue variant}
\end{figure}

\begin{figure}[H]
\centering
\includegraphics[width=0.82\textwidth]{\detokenize{results/nlp_ext/issue_transformer/nlp_issue_transformer_per_label_f1.png}}
\caption{Per-label F1 for transformer issue baseline}
\end{figure}

\clearpage
\section*{Appendix C: Syllabus Upgrade Figures}
\addcontentsline{toc}{section}{Appendix C: Syllabus Upgrade Figures}

\begin{figure}[H]
\centering
\includegraphics[width=0.82\textwidth]{\detokenize{results/nlp_ext/syllabus_upgrade/nlp_syllabus_bench_recall0.png}}
\caption{Test recall\_0 across syllabus baselines}
\end{figure}

\begin{figure}[H]
\centering
\includegraphics[width=0.82\textwidth]{\detokenize{results/nlp_ext/syllabus_upgrade/nlp_syllabus_bench_f2.png}}
\caption{Test F2\_0 across syllabus baselines}
\end{figure}

\end{document}


